\justifying
\NSFCBodyText

\subsubsection{已具备的实验条件}

申请人依托新疆大学开展研究，具备 GPU 训练环境、深度学习开发工具链以及多模态实验代码基础。SUST/RUST、Multi30k-Distant、CUTE 与 MC$^2$ 均为公开数据，写作阶段已核验来源，实验阶段按许可下载。课题组已有视觉问答与检索训练流程，可复用对比学习、评测脚本和消融记录方式。

\subsubsection{尚缺少的实验条件}

缺少现成的维语场景文字与图文融合代码；缺少 Multi30k-Distant 图像与文本的本地对齐副本；缺少 200--400 题理解协议的题面与核对记录。实验室侧的报销、年度报告与成果署名流程需在立项合同中落实。

\subsubsection{拟解决的途径}

融合代码在开源 CLIP/STR 实现上适配，不新造基础模型。图像按 Multi30k/Flickr 许可获取，不能合法使用的样本从评测中剔除并在报告中说明。理解题由研究生按测试集编写、申请人抽检。经费按申请人确定为劳务费 3.5 万元、业务费 1.5 万元，实验室实报实销；业务费对应测试、会议差旅与出版物。若劳务费不能列支，在合同阶段调整到指南允许科目。

\subsubsection{利用研究基地的计划与落实情况}

拟在立项后与“民族语言智能分析与安全治理”教育部重点实验室签订开放课题合同，按实验室要求提交年度报告与结题材料，论文第一单位使用实验室规定署名。不承诺使用实验室未对开放课题开放的内部数据或生产系统。
